% TOP_PROBLEM: {"problemId": "problem.chowla-non-vanishing-conjecture-for-dirichlet-l-functions-at-the-central-point", "canonicalTitle": "Chowla Non-Vanishing Conjecture for Dirichlet L-Functions at the Central Point", "releaseRank": 326, "primaryDomain": "number_theory_arithmetic_geometry", "primaryDomainLabel": "Number theory & arithmetic geometry", "displayStatus": "open_with_solved_subcases", "releaseStatus": "open_with_solved_subcases"}
% !TeX program = lualatex
% Definition and sources only; this file is not an LLM solution attempt.
\documentclass[11pt]{article}
\usepackage[margin=25mm]{geometry}
\usepackage{fontspec}
\setmainfont{Latin Modern Roman}
\usepackage{amsmath,amssymb,mathtools}
\usepackage{unicode-math}
\setmathfont{Latin Modern Math}
\usepackage{xurl,hyperref}
\hypersetup{colorlinks=true,urlcolor=blue}
\setcounter{secnumdepth}{0}
\setlength{\parindent}{0pt}
\setlength{\parskip}{5pt}
\setlength{\emergencystretch}{3em}
\begin{document}
\section*{\texorpdfstring{Chowla Non-Vanishing Conjecture for Dirichlet L-Functions at the Central Point}{Chowla Non-Vanishing Conjecture for Dirichlet L-Functions at the Central Point}}\label{326}
\subsection{Definitions and mathematical statement}
A Dirichlet character modulo $q$ is a homomorphism $({\mathbb{Z}}/q{\mathbb{Z}})^\times\to{\mathbb{C}}^\times$, extended periodically to integers and by zero on nonunits. It is primitive of conductor $q$ if it is not induced from a proper divisor of $q$. Initially, for $\Re s>1$,
\[
 L(s,\chi)=\sum_{n\ge1}\frac{\chi(n)}{n^s};
\]
use its analytic continuation to define the central value. Chowla's nonvanishing assertion in the selected scope is $L(\tfrac12,\chi)\ne0$ for every primitive character, real or complex, of either parity $\chi(-1)=\pm1$. The conductor-one case is the zeta function. No positivity or uniform lower bound is demanded. Primitivity matters: the statement does not simply extend the quantifier to arbitrary imprimitive characters with additional Euler factors.
\par\smallskip\textit{Formulation and concept references:} \href{https://arxiv.org/pdf/2603.22124v1}{[S1]}.

\subsection{Short English statement}
Is the central value of every primitive Dirichlet L-function nonzero, regardless of its character's parity or whether the character is real?
\subsection{Sources}
\begin{itemize}
\item[S1]Non-vanishing of Dirichlet L-functions at the central point with restricted root number; arXiv2603.22124v1,23March2026.
\url{https://arxiv.org/pdf/2603.22124v1}.
\textit{Formulation source.}
\end{itemize}
\end{document}
